%% sections/framework.tex -- Phase 7.1. Register: CLAUDE.md S9.2. The configuration framework is the
%% deliverable of the paper, given its own section (split from Discussion 2026-06-16) to foreground the
%% artifact. Centre = tab_framework; "validate locally" governs the per-axis priors. The cross-cutting
%% findings that justify the rows are discussed in Section~\ref{sec:discussion}.

\section{The Configuration Framework}
\label{sec:framework}

\begin{figure*}[b]
\centering
\resizebox{\textwidth}{!}{%
\begin{tikzpicture}[
  font=\footnotesize, >={Latex[length=2mm]},
  proc/.style={rectangle, rounded corners=2pt, draw=black!70, fill=black!6, align=center, inner sep=4pt, font=\footnotesize},
  data/.style={rectangle, rounded corners=1.5pt, draw=black!55, dashed, fill=white, align=center, font=\scriptsize, inner sep=2.5pt},
  dec/.style={diamond, aspect=1.4, draw=black!70, fill=black!9, align=center, inner sep=1pt, font=\scriptsize},
  llm/.style={rectangle, rounded corners=2pt, draw=black!85, line width=0.7pt, fill=black!14, align=center, font=\footnotesize, inner sep=4pt},
  bc/.style={rectangle, rounded corners=3pt, draw=black!60, fill=white, inner sep=4mm},
  start/.style={circle, fill=black, draw=black, minimum size=3mm, inner sep=0pt},
  edgelbl/.style={font=\scriptsize\itshape, fill=white, inner sep=1pt},
  arr/.style={->, draw=black!75, semithick},
]
% --- one block: question, answer and guidance are all the prompt; group them in one column,
%     and let the code/short decision pick which guidance the column carries ---
\node[data, text width=22mm] (q) {the question};
\node[data, text width=22mm, below=1.8mm of q] (ans) {the student answer};
\node[data, text width=22mm, below=1.8mm of ans] (rub) {a rubric};
\node[data, text width=22mm, below=1.8mm of rub] (ref) {a reference answer};
\coordinate (gmid) at ($(rub.west)!0.5!(ref.west)$);
\node[dec] (cs) at ([xshift=-16mm]gmid) {code or\\short?};
\draw[arr] (cs.east) to[out=35,in=180] node[edgelbl,above,pos=0.55,yshift=0.6mm,inner sep=0.6pt]{code} (rub.west);
\draw[arr] (cs.east) to[out=-35,in=180] node[edgelbl,below,pos=0.55,yshift=-0.6mm,inner sep=0.6pt]{short} (ref.west);
\node[draw=black!45, dotted, rounded corners=2pt, inner sep=2.2mm, fit=(q)(ans)(rub)(ref)] (ctx) {};
\node[font=\scriptsize\itshape, black!55, below=0.9mm of ctx] (ctxlbl) {prompt context};
\node[bc, fill=none, fit=(cs)(ctx)(ctxlbl)] (block) {};
\node[font=\scriptsize\itshape, anchor=south] at ([yshift=4mm]block.north) (blocklbl) {one block, ready to grade};
% stack effect: peers of the front block, offset up-right, suggesting many blocks.
% Shadows behind, then a white mask so the front block stays clean and the inputs show.
\begin{scope}[on background layer]
  \draw[bc, fill=black!5] ([xshift=3.2mm,yshift=3.2mm]block.south west) rectangle ([xshift=3.2mm,yshift=3.2mm]block.north east);
  \draw[bc, fill=black!3] ([xshift=1.6mm,yshift=1.6mm]block.south west) rectangle ([xshift=1.6mm,yshift=1.6mm]block.north east);
  \fill[white, rounded corners=3pt] (block.south west) rectangle (block.north east);
\end{scope}
% --- start, the exam, and the split, stacked at the left to save width (activity-diagram start node) ---
\node[proc, left=16mm of block, text width=24mm] (split) {Split into blocks of $n$ questions};
\node[proc, above=6mm of split] (exam) {A student's exam};
\node[start, above=6mm of exam] (start) {};
\draw[arr] (start) -- (exam);
\draw[arr] (exam) -- (split);
\draw[arr] (split.east) -- (block.west);
% --- prompt the LLM: the call applied to the assembled block (outside it) ---
\node[llm, right=16mm of block] (llm) {prompt the LLM\\[1pt]\tiny holistic $\cdot$ clean conversation};
\draw[arr] (block.east) -- (llm.west);
% reasoning is the one conditional setting of the call -> a decision above it,
% with the two outcomes placed below the diamond so they do not overlap the block
\node[dec, above=20mm of llm, xshift=9mm] (rdec) {code grader\\collapses to 0?};
\node[data, below left=5mm and 6mm of rdec] (roff) {reasoning off\\\tiny(default)};
\node[data, below right=5mm and 6mm of rdec] (ron) {reasoning on};
\draw[arr] (rdec) -- node[edgelbl,left,pos=0.45]{no} (roff.north);
\draw[arr] (rdec) -- node[edgelbl,right,pos=0.45]{yes} (ron.north);
\draw[arr] (roff.south) -- (llm.north);
\draw[arr] (ron.south) -- (llm.north);
% --- loop and the activity-diagram final node ---
\node[dec, right=14mm of llm, text width=13mm] (more) {more blocks?};
\node[circle, draw=black, minimum size=4.4mm, inner sep=0pt, right=11mm of more] (end) {};
\fill[black] (end) circle (1.1mm);
\node[font=\scriptsize\itshape, anchor=north] at ([yshift=-1mm]end.south) {exam graded};
\draw[arr] (llm) -- (more);
\draw[arr] (more) -- node[edgelbl,above]{no} (end);
\coordinate (loopbot) at ([yshift=-12mm]block.south);
\draw[arr] (more.south) |- (loopbot) -| node[edgelbl,pos=0.25,below]{yes, next block} (block.south);
\end{tikzpicture}%
}
\caption{How the recommended configuration grades an exam. The questions are split into blocks of $n$
questions (about four in the exams we studied), graded one block per call to keep the number of calls down (the
stack stands for the remaining blocks, run the same way). Each block assembles the input data the call needs, shown dashed: the question,
the student answer, and the matching guidance, a rubric for code or a reference answer for short answers.
That block is then sent to the model holistically, in a clean conversation, and the grades come back; the
process repeats block by block until the exam is graded. Reasoning is the one conditional setting: it stays
off by default and turns on only to rescue a grader that collapses to zero on code. The framework is
model-independent.}
\label{fig:frameworkdiagram}
\end{figure*}

% auto-generated by make_phase5_tables.py -- do not hand-edit
\begin{table*}[!tp]\centering\small
\caption{The configuration decision guide (Section~\ref{sec:framework}). Each row pairs a grading scenario with a recommended setting and the trade-off behind it.}\label{tab:framework}
\begin{tabular}{p{2.8cm}p{4.2cm}p{9.2cm}}
\hline
Scenario & Recommended & Trade-off in the rule \\
\hline
\raggedright Grading short answers & Reasoning off; the cheapest model & The reasoning gain is small and only sometimes significant (Qwen on Mohler and SemEval, not the other models), and reasoning off is 10--925$\times$ cheaper \\[3pt]
\hline
\raggedright Grading code, the baseline collapses to zero & Reasoning on, or switch to a discriminating model & Reasoning recovers the discrimination the collapsed baseline lost (Qwen $+0.113$ QWK on PT-CS-verified code), but at about 800$\times$ the tokens and worse run-to-run consistency, so switching model is often the better fix \\[3pt]
\hline
\raggedright Grading code, the baseline does not collapse & Reasoning off; on only if local validation justifies it & Without the collapse the reasoning gain is inconsistent and not predictable in advance, so off is the safe default, saving cost (about 94--192$\times$ the tokens) and keeping consistency \\[3pt]
\hline
\raggedright Consistency or fairness is critical & Reasoning off & Reasoning on is less reproducible, partly a length effect and partly reasoning itself \\[3pt]
\hline
\raggedright A rubric or reference answer is available & Supply it as guidance (rubric for code, reference answer for short answers) & Guidance was the one change that reliably raised agreement: it helped across datasets (RIAYN $+0.123$ QWK, PT-CS-verified code $+0.146$), at the cost of only the extra prompt tokens \\[3pt]
\hline
\raggedright The rubric has several criteria & Holistic scoring (one overall score) & On PT-CS-verified code, scoring criterion by criterion did not reliably change agreement ($-0.022$ QWK, not significant) at a higher token cost \\[3pt]
\hline
\raggedright The exam has several questions & Grade in blocks of a few questions, one call per block & On PT-CS-verified code, grouping questions into one call barely changed agreement (question-by-question to whole-exam $-0.060$ QWK, not significant), so batching them wins on cost by cutting the per-call overhead \\[3pt]
\hline
\raggedright Grading many answers in sequence & Fresh conversation per block; no history carried across the exam & A shared, accumulating conversation made grading stricter and answer order added noise (a 20-submission sub-study) \\[3pt]
\hline
\end{tabular}
\end{table*}


The study's deliverable is a configuration framework for LLM-based grading: a compact decision guide
that, for a given grading task, says which configuration to use and why. It is a small set of priors, one
per design axis, each giving a recommended default, the trade-off that justifies it, and the condition
that would call for a different choice. Table~\ref{tab:framework} sets out the priors with their
evidence, and Fig.~\ref{fig:frameworkdiagram} shows how it grades an exam. They are starting points
backed by the controlled comparisons of Section~\ref{sec:results}, not universal laws. We take them axis
by axis.

Reasoning is the axis that most repays care, and its default is off. On short answers the gain is small
while the token cost runs up to 925 times higher, so off is the clear choice. On code reasoning does
raise agreement, but on a model that already discriminates, one that spreads grades across the scale
rather than collapsing them toward zero, it is not worth the cost: GLM-5.1, the strongest code grader in
the study, gains only a small, non-significant amount from reasoning on, at several times the price, so
off stays the efficient default. Reasoning earns its keep in only one case, rescuing a weak model that collapses to scoring much of
the code at exactly zero, where it restores the discrimination the model had lost. Even there, switching
to a model that discriminates is the better fix.

The other axes resolve into simpler priors. Evaluation guidance is the one change that reliably helped, so
a trustworthy rubric or reference answer should be supplied whenever one exists. Grading scope did not
move agreement, so it is settled on efficiency: grade in blocks of a few questions rather than one at a
time, because each request carries a fixed overhead, so fewer and larger calls run faster and cost less, a
difference that compounds across a whole cohort. Scoring a rubric criterion by criterion did not
reliably change agreement with the final grade, so holistic scoring, which is also
cheaper, is the default. The grading
session should be kept clean, each answer scored on its own, because a shared conversation that accumulated
across the exam biased grades downward and answer order added noise, while clean grading was the most
reproducible regime. For the model itself, the cheapest one that discriminates is the efficient choice: an
open model with reasoning off matched the best agreement at a fraction of the cost, and the closed frontier
model never justified its much higher price.

Two design decisions give the framework its reach. First, the rows are keyed on
observable properties, whether the baseline collapses and whether a trustworthy reference exists, rather
than on model identity, so the guidance does not expire when a new model appears. Second, the guide treats consistency and cost
as first-class alongside agreement, because both moved against reasoning here and a recommendation made on
agreement alone would have pointed the wrong way. The two artifacts play complementary parts:
Fig.~\ref{fig:frameworkdiagram} shows the grading process the priors prescribe, and
Table~\ref{tab:framework} is where each choice is backed by its trade-off and the evidence behind it.

In practice the guide is a short, confident routine. Pick the cheapest model that discriminates (GLM-5.1
was the efficient all-rounder), and grade with reasoning off, a trustworthy rubric or reference where one
exists, in blocks rather than one question at a time, with holistic scores and a clean conversation per
answer.
